\begin{enumerate}
  \item 

Chaque année, le volume diminue de 10\%. Comme 10\% de \np{3000} c'est 300 (le dixième), le volume de la réserve sera égal à \mbox{$\np{3000}-300=\np{2700}$} au 1\up{er} janvier 2026, et donc \mbox{$u_1=\np{2700}$}.
  \item 

Enlever 10\% d'une quantité revient à multiplier celle-ci par \mbox{$1-\dfrac{10}{100}=\np{0.9}$}.

Donc, pour tout \mbox{$n \in \N, u_{n+1}=\np{0.9}u_n$}.
  \item 

\mbox{$(u_n)$} est donc une suite géométrique (relation de récurrence multiplicative), et sa raison est \np{0.9}.
  \item 





















On a utilisé dans la cellule \textbf{C2} la relation de récurrence, pour écrire \texttt{=0,9*B2}.
  \item 



D'après le tableau, la quantité d'eau devient inférieure à \np[~m^3]{2000} pour \mbox{$n=4$} : c'est donc à partir du 1\up{er} janvier 2029 que la quantité d’eau dans la réserve deviendra insuffisante.
\end{enumerate}  

\needspace{10\baselineskip}
